% UBT-AI-PROVENANCE-BEGIN
% schema: ubt-ai-provenance/v1
% tier: C_working
% ai_assistance: disclosed
% human_review: risk-based
% editorial_responsibility: Ing. David Jaroš
% policy: ../../../../AI_PROVENANCE.md
% notice: Working material; exhaustive human review is not claimed.
% UBT-AI-PROVENANCE-END

\documentclass[12pt]{article}
\usepackage{amsmath,amssymb}
\title{P2 – Electron Model from the $\Theta$ Field}
\author{Ing. David Jaroš}
\date{\today}
\begin{document}
\maketitle

\section*{Goal}
Demonstrate how a minimal electron-like solution can emerge from the internal structure of the unified $\Theta(q,\tau)$ field.

\section*{Approach}
\begin{itemize}
\item Use the internal spinor/tensor decomposition of $\Theta$.
\item Map quantum numbers (charge, spin, mass) to components.
\item Attempt derivation of mass term analogous to Dirac field in curved space.
\end{itemize}

\section*{Expected Outcome}
A plausible geometric derivation of electron properties as a topological excitation in $\Theta(q, \tau)$.


\section*{License}
This work is licensed under a Creative Commons Attribution 4.0 International License (CC BY 4.0).

\end{document}